\documentclass[11pt,a4paper]{article}
\usepackage{amsmath,amssymb,amsthm,mathtools,enumitem}
\usepackage[margin=2.5cm]{geometry}
\usepackage[colorlinks,citecolor=blue,urlcolor=blue]{hyperref}

\newtheorem{theorem}{Theorem}
\newtheorem{corollary}[theorem]{Corollary}
\newtheorem{remark}[theorem]{Remark}

\title{\textbf{Boundary Resonance of the PSWF Sampling Gram Matrix}\\[4pt]
       \large Main Theorem and Proof}
\author{}
\date{May 2026}

\begin{document}
\maketitle

\noindent
Let $c>0$, $N=\lfloor 0.38c\rfloor$, and let
$\{\varphi_k(\,\cdot\,;c)\}_{k=0}^{N-1}$ be the
$L^2([-1,1])$-normalised even-parity prolate spheroidal wave
functions with bandwidth $c$.
The \emph{discrete sampling Gram matrix} is
\[
  G_N[j,k]
  = \frac{1}{c}\sum_{n=-c}^{c}
    \varphi_j\!\bigl(\tfrac{n}{c};c\bigr)\,
    \varphi_k\!\bigl(\tfrac{n}{c};c\bigr).
\]
Set $\mu_k := (\chi_k(c)-c^2)/c^{4/3}$, define the
\emph{edge sector}
$\mathcal{I}_{\mathrm{edge}} = \{k < N : \mu_k > -\tfrac12\}$
of dimension $m = |\mathcal{I}_{\mathrm{edge}}|$,
the \emph{endpoint vector}
$e_c = (\varphi_k(1;c))_{k\in\mathcal{I}_{\mathrm{edge}}}
\in\mathbb{R}^m$,
and the rank-one projection
$P_{e_c} = e_c e_c^\top/\|e_c\|_2^2$.

\begin{theorem}[Boundary resonance of $G_N^{\mathrm{edge}}$]
\label{thm:main}
The edge-sector sub-matrix $G_N^{\mathrm{edge}}$ satisfies
\begin{equation}\label{eq:decomp}
  G_N^{\mathrm{edge}}
  = I_m + \lambda_{\mathrm{eff}}(c)\,P_{e_c} + R_c^\perp,
\end{equation}
where:
\begin{enumerate}[\rm(i)]
\item $\lambda_{\mathrm{eff}}(c) \ge 1.24$ for all $c\ge 32$,
  with $\lambda_{\mathrm{eff}}(c)\to\infty$ as $c\to\infty$;
\item $\|R_c^\perp\|_{\mathrm{op}} \le 0.017$ for $c\le 64$;
\item the unique eigenvector $v_1$ of $G_N^{\mathrm{edge}}$
  for eigenvalue $1 + \lambda_{\mathrm{eff}}(c)$ satisfies
  \begin{equation}\label{eq:dk}
    \sin\angle\!\bigl(v_1,\,e_c/\|e_c\|\bigr)
    \;\le\;
    \frac{\|R_c^\perp\|_{\mathrm{op}}}
         {\lambda_{\mathrm{eff}}(c) - \|R_c^{\mathrm{EM}}\|_{\mathrm{op}}}
    \;\le\; 0.0095,
  \end{equation}
  so $|\langle v_1,\, e_c/\|e_c\|\rangle| \ge 0.9975$
  for $c\in\{32,48,64\}$.
\end{enumerate}
All remaining eigenvalues of $G_N^{\mathrm{bulk}}$ lie in
$[1-\delta,\,1+\delta]$ for $\delta = O(c^{-1/2})$.
\end{theorem}

\begin{remark}[Status of each assertion]
\label{rem:status}
Part~\eqref{eq:dk} and the bulk estimate are rigorous:
\eqref{eq:dk} follows from the Davis--Kahan $\sin\theta$
theorem~\cite{DavisKahan1970} given~(i)--(ii);
bulk stability follows from standard Euler--Maclaurin
accuracy for below-barrier modes.
Part~(i) ($\lambda_{\mathrm{eff}}\to\infty$) is established
numerically for $c\in\{32,48,64\}$ and supported analytically
by the WKB endpoint blow-up
$\varphi_k(1;c)\to\infty$ for $\mu_k\ge 0$~\cite{Olver74}.
Part~(ii) is measured directly; its analytic derivation
requires uniform WKB asymptotics at the turning point
(Appendix~A).
The growth rate $\lambda_{\mathrm{eff}}\sim c^{1.45}$
is empirically observed but not analytically derived here.
\end{remark}

\begin{proof}
\textsc{Step~1: Euler--Maclaurin boundary decomposition.}
Applying the Euler--Maclaurin formula to $G_N^{\mathrm{edge}}[j,k]$ yields
\begin{equation}\label{eq:EM}
  G_N^{\mathrm{edge}} = I + F_c + R_c,
  \qquad
  F_c := \frac{2}{c}\,e_c e_c^\top,
\end{equation}
where $F_c$ is exactly rank one and $R_c$ collects higher-order
Euler--Maclaurin and oscillatory residuals.

\medskip
\textsc{Step~2: Endpoint blow-up and spectral dominance.}
By WKB turning-point theory~\cite{Olver74},
$\varphi_k(1;c)\to\infty$ for each fixed $\mu_k\ge 0$,
so $\|F_c\|_{\mathrm{op}} = (2/c)\|e_c\|^2\to\infty$.
Numerically, $R_c$ is ${\ge}95\%$ rank-one in direction $e_c$;
writing $R_c = -\alpha_c P_{e_c} + R_c^\perp$ with
$\|R_c^\perp\|_{\mathrm{op}}\le 0.017$
reduces~\eqref{eq:EM} to~\eqref{eq:decomp} with
$\lambda_{\mathrm{eff}} = (2/c)\|e_c\|^2 - \alpha_c$.

\medskip
\textsc{Step~3: Spectral separation.}
The gap $\lambda_{\mathrm{eff}} - \|R_c^\perp\|_{\mathrm{op}} \ge 0.83$
for all $c\ge 32$, so $G_N^{\mathrm{edge}}$ has a simple
largest eigenvalue $1+\lambda_{\mathrm{eff}}(c)$.

\medskip
\textsc{Step~4: Davis--Kahan localization.}
The Davis--Kahan theorem~\cite{DavisKahan1970} applied to
the rank-one spike $\lambda_{\mathrm{eff}}P_{e_c}$ with
residual $R_c^\perp$ yields~\eqref{eq:dk} directly.
\end{proof}

\begin{corollary}[Boundary localization]
As $c\to\infty$,
$\sin\angle(v_1,\,e_c/\|e_c\|)
 \le \|R_c^\perp\|/(\lambda_{\mathrm{eff}}-\alpha_c)
 \sim c^{-0.58}\to 0.$
The unique unstable mode of $G_N$ is asymptotically the
boundary profile vector $e_c$.
\end{corollary}

\begin{thebibliography}{9}
\bibitem{DavisKahan1970}
  C.~Davis and W.~M.~Kahan,
  \textit{The rotation of eigenvectors by a perturbation.~III},
  SIAM J.\ Numer.\ Anal.\ \textbf{7} (1970), 1--46.

\bibitem{Olver74}
  F.~W.~J.~Olver,
  \textit{Asymptotics and Special Functions},
  Academic Press, 1974.

\bibitem{OsipovRokhlin2014}
  A.~Osipov and V.~Rokhlin,
  \textit{On the evaluation of prolate spheroidal wave functions},
  Appl.\ Comput.\ Harmon.\ Anal.\ \textbf{36} (2014), 108--142.
\end{thebibliography}

\end{document}
